% day1_seg1_heuristic/summary.tex

\begin{frame}
  \onslide<1->{\frametitle{Mathematical Formulation of the Problem} 
  \framesubtitle{Kalman Filter Problem Summary}
  }
 \begin{itemize}
  \onslide<2->{   \item Given two models:}
    \begin{itemize}
  	\onslide<3->{	\item State Dynamics Model:  Produces predictions}
  	\onslide<4->{	\item Measurement Model:  Translates predictions into predicted measurements}
    \end{itemize}
  \onslide<5->{ 
	\item Calculate: 
	}
  \begin{itemize}
    \onslide<6->{ \item Residual, subtract the predicted measurements from 
	the actual measurements
	}
  \onslide<7->{ \item Optimal \bfit{linear} gain matrix $K(k)$ for weighting measurement residual and
      update state estimate \emph{recursively and incrementally} by adding gain-weighted residual 
	} 
	
	\end{itemize}
 \end{itemize}
\onslide<8->{\center --------------------------------------------------------------------}
\end{frame}

\begin{frame}
  \onslide<1->{\frametitle{Mathematical Formulation of the Problem} 
  \framesubtitle{Kalman Filter Problem Summary\\General, Non-Linear, Time-Invariant,
Discrete-Time Formulation}}
Given: 
 \begin{itemize}
   \onslide<2->{ \item State Dynamics Model:  $x(k)=f[x(k-1)] + w(k)$}
   \begin{itemize}
      \onslide<3->{ \item where $x$=state vector,\;\;$f$=vector-valued dynamics function,\;\;$w$=process noise vector}
   \end{itemize}
   \onslide<4->{ \item Measurement Model: $z(k)=h[x(k)]+v(k)$}
   \begin{itemize}
      \onslide<5->{ \item where $z$=measurement vector,\;\;$h$=vector-valued measurement function,\;\; $v$=measurement noise}
   \end{itemize}
  \end{itemize}
  \onslide<6->{ Calculate:}
 \begin{itemize}
  \onslide<7->{ \item Optimal gain $K(k)$ in order to update state estimate:\\
  \begin{center}
     $\stateest{k}{+} = \stateest{k}{-} + K(k)\,[\,z(k)-h(\stateest{k}{-})\,]$
	\end{center}
}

   \begin{itemize}
      \onslide<8->{ \item where $\stateest{k}{-}=f[\stateest{k-1}{+}]$=extrapolated/predicted estimate,\\ 
	  $K$=gain matrix,\;\;$h(\stateest{k}{-}])$=expected/predicted measurement}
   \end{itemize}

\end{itemize}
\onslide<9->{\center --------------------------------------------------------------------}
\end{frame}

\begin{frame}
  \onslide<1->{\frametitle{Mathematical Formulation of the Problem} 
  \framesubtitle{Kalman Filter Problem Summary:  Linear Restrictions}
  }
 \begin{itemize}
  \onslide<2->{ \item Kalman filter originally derived for linear system dynamics and linear measurement 
  model}
  \begin{itemize}
  \onslide<3->{
    \item Assume vector-valued $f$ realized as multiplication by time-invariant matrix $\Phi$ called
 \bfit{state transition matrix}
}
  \onslide<4->{
  \item Assume $h$ realized as multiplication by matrix $H(k)$
}
  \end{itemize}
  
  \onslide<5->{ \item Find:  Gain multiplier matrix $K(k)$}

\end{itemize}
\onslide<6->{\center --------------------------------------------------------------------}
\end{frame}

\begin{frame}
  \onslide<1->{\frametitle{Mathematical Formulation of the Problem} 
  \framesubtitle{Kalman Filter Problem Summary\\ Linear Discrete Formulation\\
 }
}
Given:
\begin{itemize}
  \onslide<2->{\item Discrete-Time State Dynamics Model:
  \begin{center}
  $x(k)=\Phi x(k-1)+w(k)$}
  \end{center}
  \onslide<3->{
    \begin{itemize}
		\item $x$=state vector,\;\;$\Phi$=state transition matrix,\;\; $w$=process noise vector
	\end{itemize} }
  \onslide<4->{\item Measurement Model: $z(k)=H\,x(k)+v(k)$ \\
	\begin{itemize}
		\item where $z$=measurement vector,\;\;$H$=measurement matrix (assumed constant),\;\; $v$=measurement noise vector
	\end{itemize}}
  \end{itemize}
  \onslide<5->{Calculate:}
\begin{itemize}
	  \onslide<6->{\item An optimal gain matrix $K(k)$ and new state estimate:\\
	  \begin{center}
 	$\stateest{k}{+} = \stateest{k}{-} + K(k)\,[\,z(k)-H \stateest{k}{-}\,]$ }
	  \end{center}
	\begin{itemize}
	  \onslide<7->{\item where $\stateest{k}{-}=\Phi\stateest{k-1}{+}$ extrapolates previous estimate,
			and $H \stateest{k}{-}$ is the expected/predicted measurement 
		}
	\end{itemize}
\end{itemize}
\onslide<8->{\center --------------------------------------------------------------------}
\end{frame}

\begin{frame}
  \onslide<1->{\frametitle{Mathematical Formulation of the Problem} 
\framesubtitle{Kalman Filter Problem Summary \\
Notational Reminder}
}
\begin{itemize}
  \onslide<2->{  \item Reminder:  To distinguish prior (extrapolated) estimate from updated estimate, we use 
  +/- hat notation:
} 
	\onslide<3->{  \item $\stateest{k}{-}:=\Phi \stateest{k-1}{+}$ is predicted state estimate 
	}
  \onslide<4->{  
	\item  $\stateest{k}{+}=\text{updated estimate}$
  }
  \onslide<5->{
    \item $\hat{z}(k):=H(k) \stateest{k}{-}$ sometimes used to denote expected/predicted measurment  
  }
\end{itemize}
\onslide<6->{\center --------------------------------------------------------------------}
\end{frame}

\begin{frame}
  \onslide<1->{\frametitle{Mathematical Formulation of the Problem} 
\framesubtitle{Kalman Filter Problem Summary}
}
\begin{itemize}
\onslide<2->{\item Solution should look like this:}
\onslide<3->{
\begin{equation}\label{filtersolution}
\stateest{k}{+}=\stateest{k}{-}+ K(k)\,[\,z(k)-H\stateest{k}{-}\,]
\end{equation}
}
\onslide<4-> {
  where $\stateest{k}{-}=\Phi \stateest{k-1}{+}$
}
\onslide<5-> {
  \item Looks like running average!
}
\end{itemize}
\onslide<6->{\center --------------------------------------------------------------------}
\end{frame}

\begin{frame}
  \onslide<1->{
  \onslide<1->{\frametitle{Mathematical Formulation of the Problem} 
\framesubtitle{Kalman Filter Problem Summary}
}
}
\begin{itemize}
  \onslide<2->{
\item But where are process noise $w(k)$ and measurement noise $v(k)$ used? 
}
  \onslide<3->{
\item Statistics (covariances) associated with noise processes are inputs to Kalman gain computation  
}
  \onslide<4->{
\item  Kalman gain will also be used to  
update \bfit{estimation uncertainty}, known as the \bfit{estimation error covariance} 
}
  \onslide<5->{
\item  $\Phi$ be used to predict change in estimation uncertainty over time 
}
\end{itemize}
\onslide<5->{\center --------------------------------------------------------------------}
\end{frame}
